\chapter{Feature learning and the mean-field regime of two-layer networks}\label{chap:ntkfl}

This chapter formalizes the quantitative picture of \emph{feature learning} for a two-layer network trained by
gradient descent or gradient flow, following Mei--Montanari--Nguyen (2018), Chizat--Bach (2018) and
Yang--Hu (2021). A single scaling knob $\gamma_n$ in the output interpolates between the
\emph{lazy} (neural tangent) regime $\gamma_n=1$ and the \emph{mean-field} regime $\gamma_n=\sqrt n$. The chapter proves:
the exact one-step update of the hidden-layer features and the function-space dynamics; the learning-rate condition
$\eta\asymp\gamma_n\sqrt n$ under which features move by an order-one amount; the fact that the mean-field scaling
$\gamma_n\asymp\sqrt n$, $\eta\asymp n$ is the \emph{only} one in which the features and the predictor both move at
order one; the stability of the initial residual and the freezing of the kernel below the feature-learning scale; and,
in the mean-field scaling, the rewriting of the network as an integral against the empirical measure of its neurons
together with the reduction of accelerated gradient flow to an autonomous interacting particle system with
order-one velocity.
The Lean files are in \texttt{Optimization/NTK/FeatureLearning/} (modules \texttt{Basic}, \texttt{Criterion},
\texttt{ScalingRegimes}, \texttt{Stability}, \texttt{MeanField}, \texttt{ParticleSystem}), together with the general
tools \texttt{ForMathlib/MeasureTheory/Measure/Empirical.lean},
\texttt{ForMathlib/Analysis/Asymptotics/Theta.lean} and \texttt{ForMathlib/Analysis/Calculus/Gradient/Basic.lean}.

\medskip\noindent\textbf{Notation.}
$d$ is the input dimension, $n$ the width, $m$ the number of training points. The inputs are
$x^1,\dots,x^m\in\mathbb R^d$ (a function $X:\mathrm{Fin}\,m\to\mathrm{Fin}\,d\to\mathbb R$ in Lean), the labels form
$y\in\mathbb R^m$, and $\varphi:\mathbb R\to\mathbb R$ is the activation. The hidden weights are the rows
$w_1,\dots,w_n\in\mathbb R^d$ of $W$, the readout weights $a\in\mathbb R^n$, and the flat parameter is
$\theta=\mathrm{packParams}(W,a)\in\mathbb R^{nd+n}$ (Definition~\ref{def:ntktl:packing}). The \emph{preactivation} of
neuron $i$ on sample $\alpha$ is $h_i^\alpha:=d^{-1/2}\langle w_i,x^\alpha\rangle$. The loss is
$L(\theta)=\frac1{2m}\sum_\alpha(f(x^\alpha;\theta)-y^\alpha)^2$ (Definition~\ref{def:ntkgf:outputs}).
Asymptotic statements are as $n\to\infty$ along $\mathbb N$, and $f=\Theta(g)$ means $f=O(g)$ and $g=O(f)$
(Mathlib's \texttt{IsTheta}).

\medskip\noindent\textbf{Which definitions are not introduced.}
Following the repository's convention, only two new objects are defined in this chapter:
the scaled network $f_\gamma$ (Definition~\ref{def:ntkfl:network}) and the velocity field $b$
(Definition~\ref{def:ntkfl:velocity}). The input Gram matrix $\Phi_0$, the preactivations $h_i^\alpha$, the
normalized kernel $K^{(n)}$, the one-step weights, the feature-learning criterion, the empirical measure of the neurons
and the mean-field loss $L(\rho)$ are all written out explicitly or are existing objects
(\texttt{empiricalNTKMatrix}, \texttt{mseLoss}, \texttt{MeasureTheory.empiricalMeasure}).


\section{General tools}\label{sec:ntkfl:tools}

\begin{lemma}[Asymptotic equivalence: powers and quotients]\label{lem:ntkfl:theta}
  \lean{Asymptotics.isTheta_pow_iff, Asymptotics.isTheta_iff_div_isTheta_one,
        Asymptotics.isTheta_sq_div_one_iff}
  \leanok
Let $f,g:\alpha\to\mathbb K$ ($\mathbb K$ a normed field) along a filter $l$.
(1) For $k\ne0$: $f^k=\Theta(g^k)\iff f=\Theta(g)$.
(2) If $g$ is eventually nonzero: $f=\Theta(g)\iff f/g=\Theta(1)$.
(3) If $g$ is eventually nonzero: $(f/g)^2=\Theta(1)\iff f=\Theta(g)$.
\end{lemma}

\begin{proof}
  \leanok
(1) $\Rightarrow$: if $|f|^k\le C|g|^k$ then $|f|\le C^{1/k}|g|$ (Mathlib's \texttt{IsBigO.of\_pow}), and symmetrically;
$\Leftarrow$ is \texttt{IsTheta.pow}. (2) Dividing $|f|\le C|g|$ by $|g|\neq0$ gives $|f/g|\le C$ and conversely.
(3) is (1) with $k=2$ applied to $f/g$ and $1$, followed by (2).
\end{proof}

These three facts are the only asymptotic bookkeeping needed: every statement of Sections~\ref{sec:ntkfl:criterion}
and~\ref{sec:ntkfl:regimes} is Mathlib's \texttt{IsTheta.mul}/\texttt{.div}/\texttt{.pow} plus this lemma.

\begin{definition}[Empirical measure of a finite family]\label{def:ntkfl:empirical}
  \lean{MeasureTheory.empiricalMeasure}
  \leanok
For a finite family $x=(x_i)_{i\in\iota}$ in a measurable space the \emph{empirical measure} is
$\hat\rho_x:=\frac1{|\iota|}\sum_{i\in\iota}\delta_{x_i}$.
\end{definition}

\begin{lemma}[Integrating against an empirical measure]\label{lem:ntkfl:empirical}
  \uses{def:ntkfl:empirical}
  \lean{MeasureTheory.empiricalMeasure_univ, MeasureTheory.integral_empiricalMeasure,
        MeasureTheory.lintegral_empiricalMeasure, MeasureTheory.empiricalMeasure_comp_equiv}
  \leanok
(1) If $\iota\ne\emptyset$ then $\hat\rho_x$ is a probability measure.
(2) If singletons are measurable, then for every $g:\alpha\to E$ (with $E$ a Banach space) and every $\iota$,
$\int g\,d\hat\rho_x=\frac1{|\iota|}\sum_ig(x_i)$ \emph{with no measurability or integrability hypothesis}; the
same holds for the $[0,\infty]$-valued integral.
(3) For a bijection $e:\kappa\simeq\iota$, $\hat\rho_{x\circ e}=\hat\rho_x$.
\end{lemma}

\begin{proof}
  \leanok
(1) $\hat\rho_x(\alpha)=|\iota|^{-1}\cdot|\iota|=1$. (2) A finite sum of Dirac masses on a space with measurable singletons
integrates every function: $\int g\,d\delta_{x_i}=g(x_i)$ and the integral is linear in the measure. (3) $\sum_{j\in\kappa}\delta_{x_{e(j)}}=\sum_{i\in\iota}\delta_{x_i}$
by reindexing the sum, and $|\kappa|=|\iota|$.
\end{proof}

\begin{lemma}[Scaling the output of a predictor]\label{lem:ntkfl:constMul}
  \lean{gradient_const_mul, NTK.tangentFeature_const_mul, NTK.outputJacobian_const_mul,
        NTK.empiricalNTKMatrix_const_mul}
  \leanok
Let $f$ be any model and $c\in\mathbb R$. Then $\nabla(c\,g)=c\,\nabla g$ for any scalar function $g$ on an inner product space; the tangent features of $cf$
are $c$ times those of $f$; the output Jacobian of $cf$ is $c\,J_f$; and the empirical NTK matrix of $cf$ is
$c^2K_f$.
\end{lemma}

\begin{proof}
  \uses{def:ntkgf:jacobian}
  \leanok
The gradient is linear (Mathlib only has the \texttt{fderiv} version, hence the new lemma \texttt{gradient\_const\_mul}).
The Jacobian rows are the tangent features, and $K=JJ^\top$ gives $K_{cf}=(cJ)(cJ)^\top=c^2K_f$.
\end{proof}


\section{The scaled network, single-neuron gradients and the one-step feature update}\label{sec:ntkfl:basic}

\begin{definition}[The network with scaling knob $\gamma$]\label{def:ntkfl:network}
  \uses{def:ntktl:network, def:ntktl:packing}
  \lean{NTK.featureLearningNetwork}
  \leanok
For $\gamma\in\mathbb R$ the network with scaling knob $\gamma$ is
\[
  f_\gamma(x;\theta):=\gamma^{-1}f\bigl(x/\sqrt d;\theta\bigr),\qquad\text{i.e.}\qquad
  f_\gamma(x;W,a)=\frac1{\gamma\sqrt n}\sum_{i=1}^na_i\,\varphi\Bigl(\frac{\langle w_i,x\rangle}{\sqrt d}\Bigr),
\]
where $f$ is the network of Definition~\ref{def:ntktl:network} (the Lean name is \texttt{netFromParams}, and the second formula is
\texttt{featureLearningNetwork\_packParams\_eq\_sum}). The case $\gamma=1$ is the NTK parametrization (Jacot et al.\ 2018)
and $\gamma=\sqrt n$ the mean-field parametrization, in which the prefactor is $1/n$.
\end{definition}

\begin{lemma}[Mean-field prefactor and the normalized kernel]\label{lem:ntkfl:gammaNTK}
  \uses{def:ntkfl:network, lem:ntkfl:constMul, thm:ntktl:scaledDataset}
  \lean{NTK.featureLearningNetwork_packParams, NTK.featureLearningNetwork_packParams_eq_sum,
        NTK.featureLearningNetwork_sqrt_packParams, NTK.differentiableAt_featureLearningNetwork,
        NTK.tangentFeature_featureLearningNetwork, NTK.outputJacobian_featureLearningNetwork,
        NTK.empiricalNTKMatrix_featureLearningNetwork, NTK.empiricalNTKMatrix_featureLearningNetwork_one,
        NTK.empiricalNTKMatrix_featureLearningNetwork_sqrt,
        NTK.trainingResidual_featureLearningNetwork_add}
  \leanok
Let $\tilde x^\alpha:=x^\alpha/\sqrt d$, and let $K^{(n)}(\theta)$ be the empirical NTK matrix of $f$ on the scaled dataset $(\tilde x^\alpha)$
(the \emph{normalized kernel}; by Theorem~\ref{thm:ntktl:scaledDataset} it equals the neuron average
$K^{(n),\alpha\beta}=\frac1n\sum_i[\varphi(h_i^\alpha)\varphi(h_i^\beta)+a_i^2\varphi'(h_i^\alpha)\varphi'(h_i^\beta)\Phi_0^{\alpha\beta}]$ with
$\Phi_0^{\alpha\beta}=\frac1d\langle x^\alpha,x^\beta\rangle$). Then for every $\gamma$ and $\theta$:
\begin{enumerate}
  \item (mean-field prefactor) $f_{\sqrt n}(x;W,a)=\frac1n\sum_ia_i\varphi(\langle w_i,x\rangle/\sqrt d)$;
  \item if $\varphi$ is differentiable, $\nabla_\theta f_\gamma(x;\theta)=\gamma^{-1}\nabla_\theta f(x/\sqrt d;\theta)$, and the output Jacobian of $f_\gamma$
        is $\gamma^{-1}J(\theta)$;
  \item the empirical NTK matrix of $f_\gamma$ is $\gamma^{-2}K^{(n)}(\theta)$; in particular it is $K^{(n)}$ for $\gamma=1$ and $n^{-1}K^{(n)}$ for $\gamma=\sqrt n$;
  \item the residual $r_\gamma(\theta)=f_\gamma(\theta)-y$ satisfies $r_\gamma(\theta)+y=\gamma^{-1}\bigl(r_1(\theta)+y\bigr)$, where $r_1$ is the residual of $f$ on the scaled dataset:
        the predictions are scaled by $\gamma^{-1}$, the target is not.
\end{enumerate}
\end{lemma}

\begin{proof}
  \leanok
(1) is (the formula above) with $\gamma=\sqrt n$ and $\sqrt n\sqrt n=n$. (2)--(3) are Lemma~\ref{lem:ntkfl:constMul} with $c=\gamma^{-1}$ applied to $x\mapsto f(x/\sqrt d;\theta)$, whose
NTK matrix is $K^{(n)}$ by definition of the scaled dataset. (4) is a coordinatewise computation: $f_\gamma(x^\alpha)-y^\alpha+y^\alpha=\gamma^{-1}f(\tilde x^\alpha)$.
\end{proof}

\begin{theorem}[Single-neuron gradients of the loss]\label{thm:ntkfl:gradients}
  \uses{def:ntkfl:network, lem:ntkfl:gammaNTK, thm:ntkgf:mseGradient}
  \lean{NTK.inner_tangentFeature_featureLearningNetwork, NTK.gradient_inputWeight_mseLoss,
        NTK.gradient_readout_mseLoss}
  \leanok
Let $\varphi$ be differentiable, $\gamma\in\mathbb R$, $\theta=\mathrm{packParams}(W,a)$ and write
$f^\beta=f_\gamma(x^\beta;\theta)$. For the loss $L=\mathrm{mseLoss}(f_\gamma)$ on $(X,y)$:
\begin{align*}
  \partial_{a_i}L(\theta)&=\frac1{m\,\gamma\sqrt n}\sum_{\beta=1}^m(f^\beta-y^\beta)\,\varphi(h_i^\beta),\\
  \partial_{w_{ij}}L(\theta)&=\frac1{m\,\gamma\sqrt n\sqrt d}\sum_{\beta=1}^m(f^\beta-y^\beta)\,a_i\,\varphi'(h_i^\beta)\,x^\beta_j .
\end{align*}
Moreover, for $d>0$, the Gram matrix of the tangent features is
$\langle\nabla f_\gamma^\alpha,\nabla f_\gamma^\beta\rangle=\gamma^{-2}K^{(n),\alpha\beta}$ with $K^{(n)}$ the explicit neuron average of Lemma~\ref{lem:ntkfl:gammaNTK}.
\end{theorem}

\begin{proof}
  \leanok
By Theorem~\ref{thm:ntkgf:mseGradient}, $\nabla L=\frac1m\sum_\beta(f^\beta-y^\beta)\nabla f_\gamma^\beta$, and by Lemma~\ref{lem:ntkfl:gammaNTK}(2) $\nabla f_\gamma^\beta=\gamma^{-1}\nabla f(\tilde x^\beta;\cdot)$.
Differentiating $f(\tilde x;\theta)=n^{-1/2}\sum_ia_i\varphi(\langle w_i,\tilde x\rangle)$ gives $\partial_{a_i}=n^{-1/2}\varphi(h_i)$ and $\partial_{w_{ij}}=n^{-1/2}a_i\varphi'(h_i)\,\tilde x_j$ with $\tilde x_j=x_j/\sqrt d$. The Gram statement is item~(3) of Lemma~\ref{lem:ntkfl:gammaNTK}
evaluated entrywise.
\end{proof}

\begin{theorem}[Exact one-step feature update]\label{thm:ntkfl:oneStep}
  \uses{thm:ntkfl:gradients}
  \lean{NTK.one_step_preactivation_update}
  \leanok
Let $\varphi$ be differentiable and $(W',a')$ be the parameters after one gradient step of size $\eta$,
$\mathrm{packParams}(W',a')=\mathrm{packParams}(W,a)-\eta\nabla L(W,a)$, for the loss of $f_\gamma$. Then the preactivation of hidden unit $i$ on sample $\alpha$
moves by
\[
  \Delta h_i^\alpha=d^{-1/2}\langle w_i'-w_i,x^\alpha\rangle=-\frac{\eta}{m\gamma\sqrt n}\sum_{\beta=1}^m(f^\beta-y^\beta)\,a_i\,\varphi'(h_i^\beta)\,\Phi_0^{\beta\alpha},
  \qquad\Phi_0^{\beta\alpha}=\frac1d\langle x^\beta,x^\alpha\rangle .
\]
\end{theorem}

\begin{proof}
  \leanok
By Theorem~\ref{thm:ntkfl:gradients} the row $w_i'-w_i$ equals $-\frac{\eta}{m\gamma\sqrt n\sqrt d}\sum_\beta c_{i\beta}\,x^\beta$ with
$c_{i\beta}=(f^\beta-y^\beta)a_i\varphi'(h_i^\beta)$. Taking the inner product with $x^\alpha$ and multiplying by $d^{-1/2}$ turns
$\frac1{\sqrt d}\cdot\frac1{\sqrt d}\langle x^\beta,x^\alpha\rangle$ into $\Phi_0^{\beta\alpha}$.
\end{proof}

\begin{theorem}[Function-space dynamics with the scaling knob]\label{thm:ntkfl:predictorODE}
  \uses{lem:ntkfl:gammaNTK, thm:ntkgf:outputODE}
  \lean{NTK.hasDerivAt_predictor_featureLearningNetwork}
  \leanok
Let $\varphi$ be differentiable and let $\theta(\cdot)$ follow gradient flow at rate $\eta$ for the loss of $f_\gamma$, $\theta'(t)=-\eta\nabla L(\theta(t))$. Then the training predictions
$f_\gamma(\theta(t))=(f_\gamma(x^\alpha;\theta(t)))_\alpha$ are differentiable and
\[
  \partial_tf_\gamma(\theta(t))=-\frac{\eta}{m\gamma^2}\,K^{(n)}(\theta(t))\,r_\gamma(\theta(t)).
\]
Thus the predictor moves at speed $\Theta(\eta/\gamma^2)$ times an order-one kernel, while Theorem~\ref{thm:ntkfl:oneStep}
shows that the features move at the scale $\eta/(\gamma\sqrt n)$.
\end{theorem}

\begin{proof}
  \leanok
Apply the rate-$\eta$ output ODE $\partial_tf=-\frac\eta mK_tr$ of Theorem~\ref{thm:ntkgf:outputODE} to the model $f_\gamma$, and replace its NTK matrix by $\gamma^{-2}K^{(n)}$ (Lemma~\ref{lem:ntkfl:gammaNTK}(3)).
\end{proof}

\begin{remark}[Reuse]
The unit-rate ODE chain of Chapter~\ref{chap:ntkgf} was generalized to an arbitrary rate $\eta$ for this chapter: the rate-$\eta$ theorem is the primitive,
and the unit-rate statements (\texttt{gradient\_flow\_output\_vector\_ode\_sub\_y}, \texttt{gradient\_flow\_residual\_vector\_ode},
\texttt{risk\_dissipation\_identity}) are its $\eta=1$ cases.
\end{remark}


\section{The feature-learning criterion}\label{sec:ntkfl:criterion}

\begin{remark}[Definition 4.1 (feature learning)]\label{def:ntkfl:criterion}
The hidden layer \emph{learns features} at a gradient step on sample $\alpha$ when the normalized squared displacement of the preactivations
is of order one,
\[
  \frac1n\|\Delta h^\alpha\|^2=\frac1n\sum_{i=1}^n(\Delta h_i^\alpha)^2=\Theta(1)\qquad(n\to\infty).
\]
No separate definition is made in Lean: the criterion is Mathlib's \texttt{IsTheta} along \texttt{atTop} applied to the sequence $n\mapsto n^{-1}\sum_{i<n}(\Delta h_i^\alpha)^2$ and the constant $1$,
with $\Delta h$ the explicit expression of Theorem~\ref{thm:ntkfl:oneStep}.
\end{remark}

\begin{theorem}[Exact scaling of the normalized displacement]\label{thm:ntkfl:displacement}
  \uses{lem:ntkfl:theta}
  \lean{NTK.normalized_sq_displacement_isTheta}
  \leanok
Suppose $\Delta h_i^{(n)}=c_nS_i^{(n)}$ for all $n,i$ and the neuron average $\frac1n\sum_i(S_i^{(n)})^2=\Theta(1)$ (typical coordinates are of order one).
Then $\frac1n\sum_i(\Delta h_i^{(n)})^2=\Theta(c_n^2)$.
\end{theorem}

\begin{proof}
  \leanok
$\frac1n\sum_i(c_nS_i)^2=c_n^2\cdot\frac1n\sum_iS_i^2$, and a product of $c_n^2$ with a $\Theta(1)$ sequence is $\Theta(c_n^2)$.
\end{proof}

\begin{theorem}[Fundamental learning-rate condition]\label{thm:ntkfl:criterion}
  \uses{thm:ntkfl:oneStep, thm:ntkfl:displacement, lem:ntkfl:theta}
  \lean{NTK.featureLearning_criterion_iff_learningRate, NTK.featureLearning_iff_learningRate_of_oneStep}
  \leanok
(1) For positive sequences $\gamma_n$ and any $\eta_n$:
$\bigl(\eta_n/(\gamma_n\sqrt n)\bigr)^2=\Theta(1)\iff\eta_n=\Theta(\gamma_n\sqrt n)$.
(2) Let $\varphi$ be differentiable, $m>0$, and for each $n$ let $(W_n',a_n')$ be obtained from $(W_n,a_n)$ by one gradient step of size $\eta_n$ on the loss of $f_{\gamma_n}$ (fixed data $X,y$).
Fix a sample $\alpha$ and write $S_i^{(n)}:=\sum_\beta(f^\beta-y^\beta)\,a_i\,\varphi'(h_i^\beta)\,\Phi_0^{\beta\alpha}$. If $\frac1n\sum_i(S_i^{(n)})^2=\Theta(1)$, then
\[
  \frac1n\sum_{i=1}^n(\Delta h_i^\alpha)^2=\Theta(1)\iff\eta_n=\Theta(\gamma_n\sqrt n).
\]
\end{theorem}

\begin{proof}
  \leanok
(1) is Lemma~\ref{lem:ntkfl:theta}(3) with $f=\eta$, $g=\gamma\sqrt n$ (nonzero since $\gamma_n>0$). (2) By Theorem~\ref{thm:ntkfl:oneStep}, $\Delta h_i=c_nS_i$ with $c_n=-\eta_n/(m\gamma_n\sqrt n)$.
Theorem~\ref{thm:ntkfl:displacement} gives $\frac1n\sum_i\Delta h_i^2=\Theta(c_n^2)$ and $c_n^2=m^{-2}(\eta_n/(\gamma_n\sqrt n))^2$ is $\Theta$ of $(\eta_n/(\gamma_n\sqrt n))^2$ because $m^{-2}\ne0$ is a constant. Combine with (1).
\end{proof}


\section{The scaling regimes and the unique balance}\label{sec:ntkfl:regimes}

Two quantities govern the dynamics (Theorems~\ref{thm:ntkfl:oneStep} and~\ref{thm:ntkfl:predictorODE}): the \emph{feature motion} $(\eta/(\gamma\sqrt n))^2$ (the normalized
squared displacement of the features, up to the constant $m^{-2}$) and the \emph{function motion} $\eta/\gamma^2$ (the speed of the predictor).

\begin{theorem}[Master scaling lemma]\label{thm:ntkfl:master}
  \lean{NTK.scales_isTheta}
  \leanok
If $\gamma=\Theta(\gamma')$ and $\eta=\Theta(\eta')$ then $(\eta/(\gamma\sqrt n))^2=\Theta\bigl((\eta'/(\gamma'\sqrt n))^2\bigr)$ and $\eta/\gamma^2=\Theta(\eta'/\gamma'^2)$.
\end{theorem}

\begin{proof}
  \leanok
Quotients, products with $\sqrt n$, and powers preserve $\Theta$ (\texttt{IsTheta.div}, \texttt{.mul}, \texttt{.pow}).
\end{proof}

\begin{theorem}[The three regimes]\label{thm:ntkfl:regimes}
  \uses{thm:ntkfl:master, thm:ntkfl:criterion}
  \lean{NTK.scaling_lazy_ntk_dynamics, NTK.isLittleO_sqrt_of_lazy_learningRate,
        NTK.scaling_naive_large_eta_instability, NTK.scaling_mean_field_isTheta_one,
        NTK.scaling_mean_field_muP_balance}
  \leanok
\begin{enumerate}
  \item (Lazy / NTK: $\gamma=1$, $\eta=\Theta(1)$.) Feature motion $=\Theta(n^{-1})\to0$ and function motion $=\Theta(1)$; moreover $\eta=O(1)$ implies $\eta=o(\sqrt n)$, i.e.\ $\eta$ is below the feature-learning scale.
  \item (Naive large rate: $\gamma=1$, $\eta=\Theta(\sqrt n)$.) Feature motion $=\Theta(1)$ (features are learned), but function motion $=\Theta(\sqrt n)\to\infty$: the predictor dynamics blow up.
  \item (Mean field / $\mu$P: $\gamma=\sqrt n$, $\eta=\Theta(n)$, e.g.\ $\eta=\eta_0n$ with $\eta_0\ne0$.) Both feature motion and function motion are $\Theta(1)$.
\end{enumerate}
\end{theorem}

\begin{proof}
  \leanok
Each case is Theorem~\ref{thm:ntkfl:master} against a reference pair $(\gamma',\eta')$. (1) With $\gamma=1$: $(1/\sqrt n)^2=1/n$, and $1/n\to0$. (2) With $\gamma=1$, $\eta=\Theta(\sqrt n)$, Theorem~\ref{thm:ntkfl:criterion}(1) gives feature motion $\Theta(1)$, and $\eta/1=\Theta(\sqrt n)\to\infty$ in norm. (3) With $\gamma=\sqrt n$:
$\sqrt n\sqrt n=n$ and $\sqrt n^2=n$, so both scales are $(\eta/n)^2$ and $\eta/n$, which are $\Theta(1)$ exactly when $\eta=\Theta(n)$.
The statement $\eta=O(1)\Rightarrow\eta=o(\sqrt n)$ is Mathlib's \texttt{IsBigO.trans\_isLittleO} with $1=o(\sqrt n)$.
\end{proof}

\begin{theorem}[Unique balanced scaling]\label{thm:ntkfl:balance}
  \uses{thm:ntkfl:criterion, thm:ntkfl:regimes, lem:ntkfl:theta}
  \lean{NTK.featureMotion_and_functionMotion_isTheta_one_iff, NTK.unique_scaling_balance_solution,
        NTK.featureMotion_and_functionMotion_isTheta_one_iff_meanField}
  \leanok
Let $\gamma_n$ be eventually nonzero. Then
\[
  \Bigl(\bigl(\tfrac{\eta_n}{\gamma_n\sqrt n}\bigr)^2=\Theta(1)\ \wedge\ \tfrac{\eta_n}{\gamma_n^2}=\Theta(1)\Bigr)
  \iff\bigl(\eta=\Theta(\gamma\sqrt n)\ \wedge\ \eta=\Theta(\gamma^2)\bigr)
  \iff\bigl(\gamma=\Theta(\sqrt n)\ \wedge\ \eta=\Theta(n)\bigr).
\]
That is, the mean-field scaling $\gamma\asymp\sqrt n$, $\eta\asymp n$ is the \emph{only} one in which the features and the predictor both move at order one.
\end{theorem}

\begin{proof}
  \leanok
The first equivalence is Lemma~\ref{lem:ntkfl:theta}(3) for the feature motion and (2) for the function motion ($\eta/\gamma^2=\Theta(1)\iff\eta=\Theta(\gamma^2)$).
For the second, ``$\Leftarrow$'': $\gamma=\Theta(\sqrt n)$ gives $\gamma\sqrt n=\Theta(n)$ and $\gamma^2=\Theta(n)$, so $\eta=\Theta(n)$ is $\Theta$ of both. ``$\Rightarrow$'': from
$\eta=\Theta(\gamma\sqrt n)$ and $\eta=\Theta(\gamma^2)$ we get $\gamma\sqrt n=\Theta(\gamma^2)$; dividing by the eventually nonzero $\gamma$ yields $\sqrt n=\Theta(\gamma)$, and then $\eta=\Theta(\gamma\sqrt n)=\Theta(\sqrt n\sqrt n)=\Theta(n)$.
\end{proof}


\section{Stability: the initial residual and the kernel drift}\label{sec:ntkfl:stability}

\begin{theorem}[Initial residual stability]\label{thm:ntkfl:initialResidual}
  \uses{lem:ntkfl:gammaNTK, thm:ntktl:initialWeakLimit}
  \lean{NTK.tendsto_initial_trainingResidual_add_target}
  \leanok
Let $\varphi$ be measurable with $w\mapsto\varphi(\langle w,\tilde x^\alpha\rangle)$ in $L^2(\mathcal N(0,I_d))$ for each $\alpha$ (with $\tilde x^\alpha=x^\alpha/\sqrt d$), and let $\gamma_n\to\infty$ (for example the mean-field value $\gamma_n=\sqrt n$).
Under the Gaussian initialization $\mu_{n,d}$ (Definition~\ref{def:ntktl:init}), for every $\varepsilon>0$,
\[
  \mu_{n,d}\bigl(\|r_{\gamma_n}(0)+y\|\ge\varepsilon\bigr)\longrightarrow0\qquad(n\to\infty).
\]
In words: the initial prediction of $f_{\gamma_n}$ is negligible, so the initial residual is $-y+o_P(1)$.
\end{theorem}

\begin{proof}
  \leanok
By Lemma~\ref{lem:ntkfl:gammaNTK}(4), $r_{\gamma_n}(0)+y=\gamma_n^{-1}(r_1(0)+y)$ with $r_1$ the residual of $f$ on the scaled dataset. By Theorem~\ref{thm:ntktl:initialWeakLimit} (the part giving a width-independent radius), for every $\delta>0$ there is $R$ with
$\mu_{n,d}(\|r_1(0)\|>R)\le\delta$ for all $n$. On the complement, $\|r_1(0)+y\|\le R+\|y\|$, so $\|r_{\gamma_n}(0)+y\|\le(R+\|y\|)/\gamma_n<\varepsilon$ as soon as $\gamma_n>(R+\|y\|)/\varepsilon$, which holds eventually.
Hence the probability is at most $\delta$ for all large $n$.
\end{proof}

\begin{theorem}[The kernel drift of one step is controlled by the feature motion]\label{thm:ntkfl:kernelDrift}
  \uses{thm:ntkfl:gradients, thm:ntktl:jacobianLipschitz, thm:ntkgf:kernelLipschitz, thm:ntkgf:mseGradient}
  \lean{NTK.norm_oneStep_displacement_le, NTK.kernel_drift_oneStep_le}
  \leanok
Let $n>0$, $C_1,C_2,R\ge0$, and let $\varphi$ be differentiable and $C_1$-Lipschitz, $|\varphi'|\le C_1$ and $\varphi'$ $C_2$-Lipschitz; let $\gamma>0$, $\eta\in\mathbb R$, and $\theta_1=\theta_0-\eta\nabla L(\theta_0)$ be one step on the loss of $f_\gamma$. Write $J$ for the Jacobian of $f$ on the scaled dataset and
$K^{(n)}=JJ^\top$. Assume $\|J(\theta_0)\|,\|J(\theta_1)\|\le M$, $\|r_\gamma(\theta_0)\|\le\rho$, the readouts of $\theta_1$ are bounded by $R$, and $K_c\ge\bigl(\sum_\alpha(2R^2C_2^2\|\tilde x^\alpha\|^4+3C_1^2\|\tilde x^\alpha\|^2)\bigr)^{1/2}$. Then (with Frobenius norms)
\begin{align*}
  \|\theta_1-\theta_0\|&\le\frac{|\eta|}{\gamma}\cdot\frac{M\rho}{m},\\
  \|K^{(n)}(\theta_1)-K^{(n)}(\theta_0)\|&\le\frac{2M^2\rho K_c}{m}\cdot\Bigl|\frac{\eta}{\gamma\sqrt n}\Bigr|\qquad(n>0).
\end{align*}
The last factor is the square root of the feature motion of Section~\ref{sec:ntkfl:regimes}.
\end{theorem}

\begin{proof}
  \leanok
Displacement: $\theta_1-\theta_0=-\eta\nabla L(\theta_0)$ and, by Theorem~\ref{thm:ntkgf:mseGradient}, $\|\nabla L\|\le m^{-1}\|J_\gamma\|\,\|r_\gamma\|$ with $J_\gamma=\gamma^{-1}J$ (Lemma~\ref{lem:ntkfl:gammaNTK}(2)), giving $\|\nabla L\|\le\gamma^{-1}m^{-1}M\rho$.
Kernel: by Theorem~\ref{thm:ntktl:jacobianLipschitz}, $\|J(\theta_1)-J(\theta_0)\|\le(K_c/\sqrt n)\|\theta_1-\theta_0\|$, and by Theorem~\ref{thm:ntkgf:kernelLipschitz}, $\|K(\theta_1)-K(\theta_0)\|\le2M\|J(\theta_1)-J(\theta_0)\|$. Chaining,
$\|K(\theta_1)-K(\theta_0)\|\le2M\cdot\frac{K_c}{\sqrt n}\cdot\frac{|\eta|}{\gamma}\frac{M\rho}m=\frac{2M^2\rho K_c}{m}\bigl|\frac{\eta}{\gamma\sqrt n}\bigr|$.
\end{proof}

\begin{theorem}[Kernel freeze below the feature-learning scale]\label{thm:ntkfl:freeze}
  \uses{thm:ntkfl:kernelDrift, thm:ntkfl:regimes}
  \lean{NTK.tendsto_kernel_drift_of_isLittleO}
  \leanok
Under the hypotheses of Theorem~\ref{thm:ntkfl:kernelDrift} with constants $C_1,C_2,R,M,\rho,K_c$ uniform in the width (the bounds on $J$ and $r$ holding eventually in $n$) and $\gamma_n>0$, if
$\eta_n=o(\gamma_n\sqrt n)$, then $\|K^{(n)}(\theta_1^{(n)})-K^{(n)}(\theta_0^{(n)})\|\to0$. In particular this holds in the lazy regime $\gamma=1$, $\eta=O(1)$ (Theorem~\ref{thm:ntkfl:regimes}(1)).
\end{theorem}

\begin{proof}
  \leanok
$|\eta_n/(\gamma_n\sqrt n)|\to0$ because $\eta_n=o(\gamma_n\sqrt n)$ (\texttt{IsLittleO.tendsto\_div\_nhds\_zero}); the drift is squeezed between $0$ and $\frac{2M^2\rho K_c}{m}|\eta_n/(\gamma_n\sqrt n)|$ by Theorem~\ref{thm:ntkfl:kernelDrift}.
\end{proof}

\begin{remark}[What is not formalized: the converse of the freeze]\label{rmk:ntkfl:noConverse}
The informal dichotomy ``kernel frozen iff features not learned'' has two halves. The half above (below the feature-learning scale the kernel is frozen, with an explicit rate) is proved. The converse (feature learning, $\eta=\Theta(\gamma\sqrt n)$, forces $\|K_1-K_0\|\not\to0$) is \emph{not} formalized and does not follow from the criterion alone: order-one motion of every feature does not force a change of the Gram matrix, since neuron relabelling, sign flips for an even activation, or an orthogonal mixing of hidden units can move features while leaving $K^{(n)}$ unchanged. A lower bound therefore needs a non-degeneracy hypothesis on $\varphi$ and the data, whose verification at Gaussian initialization would require a law of large numbers over neurons for the second-order kernel increment (the first-order terms cancel by the random signs of $a_i$) and the exclusion of coincidental cancellation. Likewise, Theorem~\ref{thm:ntkfl:initialResidual} is the convergence-in-probability form; the explicit $O_P(n^{-1/2})$ rate for the initial prediction is not formalized.
\end{remark}


\section{The mean-field description: the network as an integral over its neurons}\label{sec:ntkfl:meanfield}

In the mean-field scaling $\gamma=\sqrt n$ the network depends on its neurons $p_i=(a_i,w_i)\in\Theta:=\mathbb R\times\mathbb R^d$ only through their empirical measure
$\rho^{(n)}:=\frac1n\sum_i\delta_{p_i}$ (Definition~\ref{def:ntkfl:empirical}). For a general finite measure $\rho$ on $\Theta$ the \emph{mean-field predictor} is
\[
  f_\rho(x):=\int_\Theta a\,\varphi\Bigl(\frac{\langle w,x\rangle}{\sqrt d}\Bigr)\,\rho(da\,dw).
\]
No Lean definition is made: the integral $\int p_1\,\varphi(d^{-1/2}\langle p_2,x\rangle)\,d\rho(p)$ is written out in each statement, and the mean-field loss
$L(\rho)=\frac1{2m}\sum_\alpha(f_\rho(x^\alpha)-y^\alpha)^2$ is the existing \texttt{mseLoss} applied to the function $(x,\rho)\mapsto f_\rho(x)$, which is possible because \texttt{mseLoss} and \texttt{trainingResidual} are polymorphic in the parameter space (here $\Theta=$ the space of measures on $\mathbb R\times\mathbb R^d$). (As a Bochner integral $f_\rho$ is $0$ where its integrand is not integrable; this never matters for the empirical measures treated here.)

\begin{theorem}[Integral representation of the width-$n$ network (Definition 4.2)]\label{thm:ntkfl:integralRep}
  \uses{def:ntkfl:empirical, lem:ntkfl:empirical, lem:ntkfl:gammaNTK}
  \lean{NTK.integral_empiricalMeasure_neuron,
        NTK.integral_empiricalMeasure_neuron_eq_featureLearningNetwork}
  \leanok
For every $n,d$, activation $\varphi$, readout weights $a$, hidden weights $W$ and input $x$ (no hypothesis on $\varphi$, and $n=0$ allowed),
\[
  f_{\rho^{(n)}}(x)=\frac1n\sum_{i=1}^na_i\,\varphi\Bigl(\frac{\langle w_i,x\rangle}{\sqrt d}\Bigr)=f_{\sqrt n}\bigl(x;\mathrm{packParams}(W,a)\bigr).
\]
\end{theorem}

\begin{proof}
  \leanok
The first equality is Lemma~\ref{lem:ntkfl:empirical}(2) applied to $g(p)=p_1\,\varphi(d^{-1/2}\langle p_2,x\rangle)$ and the family $p_i=(a_i,w_i)$, which needs no integrability hypothesis; for $n=0$ both sides are $0$. The second is Lemma~\ref{lem:ntkfl:gammaNTK}(1).
\end{proof}

\begin{theorem}[The mean-field loss at the empirical measure is the network loss]\label{thm:ntkfl:lossRep}
  \uses{thm:ntkfl:integralRep}
  \lean{NTK.trainingResidual_integral_empiricalMeasure, NTK.mseLoss_integral_empiricalMeasure}
  \leanok
With $\rho^{(n)}=\hat\rho_{(a_i,w_i)}$, the training residual and the empirical MSE loss of $\rho\mapsto(f_\rho(x^\alpha))_\alpha$ at $\rho^{(n)}$ are those of the width-$n$ mean-field network $f_{\sqrt n}$ at $\mathrm{packParams}(W,a)$:
$r(\rho^{(n)})=r_{\sqrt n}(W,a)$ and $L(\rho^{(n)})=L(W,a)$.
\end{theorem}

\begin{proof}
  \leanok
Both sides are built from the vector of predictions on the data, which agree coordinatewise by Theorem~\ref{thm:ntkfl:integralRep}; the loss is $(2m)^{-1}\|r\|^2$ of the residual.
\end{proof}


\section{Accelerated gradient flow as an interacting particle system}\label{sec:ntkfl:particles}

\begin{definition}[The mean-field velocity field]\label{def:ntkfl:velocity}
  \uses{thm:ntkfl:integralRep}
  \lean{NTK.meanFieldVelocityField}
  \leanok
For $\eta_0\in\mathbb R$, a particle $p=(a,w)\in\mathbb R\times\mathbb R^d$ and a measure $\rho$ on $\mathbb R\times\mathbb R^d$ let $h^\beta:=d^{-1/2}\langle w,x^\beta\rangle$ and define $b(p,\rho)=(b_a,b_w)$ by
\[
  b_a(p,\rho)=-\frac{\eta_0}{m}\sum_{\beta=1}^m\bigl(f_\rho(x^\beta)-y^\beta\bigr)\varphi(h^\beta),\qquad
  b_w(p,\rho)=-\frac{\eta_0}{m\sqrt d}\sum_{\beta=1}^m\bigl(f_\rho(x^\beta)-y^\beta\bigr)\,a\,\varphi'(h^\beta)\,x^\beta .
\]
\end{definition}

\begin{theorem}[Cancellation of the width factor]\label{thm:ntkfl:cancellation}
  \uses{def:ntkfl:velocity, thm:ntkfl:gradients, thm:ntkfl:integralRep}
  \lean{NTK.width_factor_cancellation}
  \leanok
Let $\varphi$ be differentiable, $L$ the loss of the mean-field network $f_{\sqrt n}$ and $\theta=\mathrm{packParams}(W,a)$. For every neuron $i$ (so $n>0$) and every $m,d$,
\[
  b\bigl((a_i,w_i),\rho^{(n)}\bigr)=\bigl(-\eta_0n\,\partial_{a_i}L(\theta),\ -\eta_0n\,\nabla_{w_i}L(\theta)\bigr).
\]
That is, multiplying the single-neuron loss gradients (which carry the width factor $1/n$) by the accelerated rate $\eta_0n$ leaves an order-one velocity depending on the other neurons only through $\rho^{(n)}$.
\end{theorem}

\begin{proof}
  \leanok
Theorem~\ref{thm:ntkfl:gradients} with $\gamma=\sqrt n$ gives $\partial_{a_i}L=\frac1{m\,n}\sum_\beta(f^\beta-y^\beta)\varphi(h_i^\beta)$ and $\partial_{w_i}L=\frac1{m\,n\sqrt d}\sum_\beta(f^\beta-y^\beta)a_i\varphi'(h_i^\beta)x^\beta$, because $\sqrt n\sqrt n=n$. By Theorem~\ref{thm:ntkfl:integralRep} each $f^\beta=f_{\sqrt n}(x^\beta)$ equals $f_{\rho^{(n)}}(x^\beta)$. Multiplying by $-\eta_0n$ cancels $n^{-1}$ and gives $b_a$, $b_w$ of Definition~\ref{def:ntkfl:velocity}.
\end{proof}

\begin{theorem}[The interacting particle system]\label{thm:ntkfl:particleODE}
  \uses{thm:ntkfl:cancellation, thm:ntkgf:outputODE}
  \lean{NTK.interacting_particle_system_ode}
  \leanok
Let $\varphi$ be differentiable and let $t\mapsto\theta(t)=\mathrm{packParams}(W(t),a(t))$ follow gradient flow at the accelerated rate $\eta_0n$ for the loss of $f_{\sqrt n}$: $\theta'(t)=-\eta_0n\nabla L(\theta(t))$.
Then for every neuron $i$ the particle $t\mapsto(a_i(t),w_i(t))$ is differentiable and
\[
  \frac{d}{dt}\bigl(a_i(t),w_i(t)\bigr)=b\Bigl(\bigl(a_i(t),w_i(t)\bigr),\ \rho_t^{(n)}\Bigr),\qquad
  \rho_t^{(n)}=\frac1n\sum_{j=1}^n\delta_{(a_j(t),w_j(t))}.
\]
\end{theorem}

\begin{proof}
  \leanok
A curve in Euclidean space is differentiable with derivative $v$ iff each coordinate is, with derivative $v_k$ (\texttt{hasDerivAt\_euclideanSpace}). The coordinates of $\theta(t)$ at the indices of $a_i$ and $w_{ij}$ are $a_i(t)$ and $w_{ij}(t)$ (Definition~\ref{def:ntktl:packing}), and their derivatives are $-\eta_0n\,\partial_{a_i}L$ and $-\eta_0n\,\partial_{w_{ij}}L$ at $\theta(t)$. Theorem~\ref{thm:ntkfl:cancellation} identifies these with $b$ at $((a_i(t),w_i(t)),\rho_t^{(n)})$; the pair of components is a differentiable curve into $\mathbb R\times\mathbb R^d$.
\end{proof}

\begin{theorem}[Structure of the coupling]\label{thm:ntkfl:coupling}
  \uses{def:ntkfl:velocity, lem:ntkfl:empirical}
  \lean{NTK.meanFieldVelocityField_congr, NTK.mean_field_coupling_structure}
  \leanok
(1) If two measures $\rho,\rho'$ have the same predictions $f_\rho(x^\beta)=f_{\rho'}(x^\beta)$ on every data point then $b(p,\rho)=b(p,\rho')$ for every particle $p$:
a particle feels the others \emph{only} through the scalar predictions on the data.
(2) For any permutation $\sigma$ of the $n$ neurons, $b\bigl(p,\hat\rho_{(a_{\sigma(j)},w_{\sigma(j)})}\bigr)=b\bigl(p,\hat\rho_{(a_j,w_j)}\bigr)$: the coupling is all-to-all with the permutation-symmetric weight $1/n$ carried by the empirical measure.
\end{theorem}

\begin{proof}
  \leanok
(1) $\rho$ enters $b$ only through $f_\rho(x^\beta)$. (2) is Lemma~\ref{lem:ntkfl:empirical}(3).
\end{proof}

\begin{remark}[What is and is not formalized in the mean-field part]
The statements above are exact identities at each finite width $n$. Not formalized: the existence of the gradient flow itself (Theorem~\ref{thm:ntkfl:particleODE} takes the flow as a hypothesis), and the mean-field \emph{limit} $n\to\infty$ (propagation of chaos / convergence of $\rho_t^{(n)}$ to a solution of the limiting continuity equation), which is the content of Mei--Montanari--Nguyen and Chizat--Bach beyond the structural reduction proved here.
\end{remark}
